\subsubsection{Small game tables (Kayles SG values)}

\paragraph{Idea intuitiva.}
Cuando el juego tiene \textbf{pocas} posiciones conocidas, conviene \textbf{hardcodear} la tabla de Grundys para responder en $O(1)$ por query. Kayles es un juego clasico donde tiras un bolo y tumbas uno o dos adyacentes; sus Grundys para $n \le 71$ son bien conocidos y estan precomputados. La razon de precomputar: el Grundy de Kayles satisface una recurrencia compleja y no se puede calcular con formula cerrada simple.

\paragraph{Propiedades clave (invariantes).}
\begin{itemize}\setlength\itemsep{0pt}
	\item Solo util para Kayles; otros juegos necesitan su propia tabla.
	\item La tabla se obtuvo experimentalmente (o por recurrencia) y se almacena literal en el codigo.
	\item Los Grundys tienen \textbf{periodo} tras cierto $n$ (e.g., Kayles tiene periodo $12$ tras $n = 72$).
\end{itemize}

\paragraph{Codigo clave.}
Si el problema tiene $n$ en $[1, N]$ con $N$ conocido, precomputar la tabla:
\begin{verbatim}
vector<int> grundy(N + 1, 0);
for (int n = 1; n <= N; ++n) {
    set<int> s;
    for (int k = 0; k < n; ++k) {
        // un bolo en k: kayles se parte en [0, k-1] y [k+1, n-1]
        s.insert(grundy[k] ^ grundy[n-1-k]);
        if (k+1 < n) s.insert(grundy[k] ^ grundy[n-2-k]);
    }
    int g = 0; while (s.count(g)) ++g;
    grundy[n] = g;
}
\end{verbatim}

\paragraph{Complejidad.}
\begin{itemize}\setlength\itemsep{0pt}
	\item $O(1)$ por query (con tabla hardcodeada).
	\item Precompute $O(N^2)$.
\end{itemize}

\paragraph{Donde tocar para modificar.}
\begin{itemize}\setlength\itemsep{0pt}
	\item \textbf{Otro juego}: aniadir otra tabla literal (Nim, Dawson's Kayles, Grundy's Game, etc.).
	\item \textbf{Tabla generada por codigo}: aniadir un precompute al inicio que llene la tabla con el template Sprague-Grundy.
	\item \textbf{Usar periodicidad}: si el Grundy tiene periodo $P$ a partir de $N_0$, aniadir \texttt{int g = (n >= N0) ? tabla[(n - N0) \% P] : tabla[n];}.
\end{itemize}

\paragraph{Trampas y errores frecuentes.}
\begin{itemize}\setlength\itemsep{0pt}
	\item Si el problema da $n$ fuera de rango, aniadir bounds check antes de indexar.
	\item La tabla es especifica de la variante de Kayles (original, Dawson's, etc.); verificar cual es.
	\item En algunos juegos (Wythoff), la estrategia no es Grundy sino un par de moves explicitos.
\end{itemize}

\paragraph{Explicacion linea por linea.}
\begin{itemize}\setlength\itemsep{0pt}
	\item \verb|#include <bits/stdc++.h>| -- incluye todas las bibliotecas estandar.
	\item \verb|using namespace std;| -- evita prefijar \verb|std::|.
	\item \verb|vector<int> kaylesSG = \{...\}| -- tabla hardcodeada con los Grundys de Kayles para $n = 0, \dots, 71$ (precalculada fuera de linea).
	\item Los valores en la tabla siguen la recurrencia SG de Kayles: cada $n$ tiene un Grundy unico que codifica las estrategias ganadoras/perdedoras.
	\item \verb|cout << "SG Kayles(7) = " << kaylesSG[7] << "\textbackslash n";| -- consulta en $O(1)$: el Grundy de Kayles con 7 bolos es $3$.
\end{itemize}
Para entender la teoria de la tabla: el Grundy $g(n)$ de Kayles se calcula como $g(n) = \text{mex}\{g(k) \oplus g(n-1-k), g(k) \oplus g(n-2-k)\}$ para todos los $k$ validos; los valores en la tabla son la solucion de esta recurrencia (eventualmente periodica con periodo 12).

\paragraph{Codigo.}
Ver \texttt{cpp.tex}, seccion \textit{...}.
